\subsection{Complete census of cylinders and selection of the terminal boundary}
\label{act21:censo:terminal-completo}

The terminal readout jointly uses two kinds of information: compatibility
of the regional cylinders and an incidence profile with labelled channels and
positions. This section constructs the complete catalogues,
enumerates their triples, and proves the effect of each reader condition.
The starting point consists of the regional outputs already produced in
Section~\ref{sec:generacion}, not values introduced to define
the APP--TRIT--TPK transport. Reduction to a matrix belongs to a subsequent
readout of these outputs and preserves the order of the three channels.

The result is a census theorem with an explicit starting structure.
The profile of Gram products, weights, distances, and marginals is declared as a datum of the
terminal reader. The exhaustiveness of the enumeration proves what
this profile selects on the specified catalogues; it does not, by itself, prove that the
profile is a universal law or that it can be reconstructed from the unique triple
that ultimately satisfies it. This distinction allows the complete finite
proof to be incorporated without reversing its genealogy.

\subsubsection{From regional outputs to ternary catalogues}

Let \(x_\pi,x_e,x_\varphi\in[0,1)\) be the fractional parts of the
three regional coordinates already constructed. We use their expressed
cylinders of one thousand decimal places,
\[
 J_\chi^{(1000)}=
 \left[\frac{d_\chi}{10^{1000}},
       \frac{d_\chi+1}{10^{1000}}\right),
 \qquad \chi\in\{\pi,e,\varphi\}.
\]
The integers \(d_\chi\) are finite output readouts. A length of one thousand
is sufficient for the integers that follow, as verified from the
two bounds of each interval; no infinite depth is assumed.
We define
\begin{equation}
 p=30,\quad r=1050,\quad L=p+r=1080,\quad k=171,
 \qquad D=3^L,\quad Q=1000^k,
 \label{act21:censo:parametros}
\end{equation}
and
\[
 P_\chi=\lfloor3^{30}x_\chi\rfloor,\qquad
 T_\chi=\lfloor Qx_\chi\rfloor.
\]
Here \(k\) counts decimal triads and does not denote the dodecaphasic register
\(K\). We have \(1000^{171}\leq3^{1080}<1000^{172}\).

\begin{lemma}[Stable extraction of an integer]
\label{act21:censo:estabilidad}
For \(d,m,h\in\mathbb Z\), with \(m,h>0\), the value of
\(\lfloor mx\rfloor\) is constant on \([d/h,(d+1)/h)\) if and only if
\begin{equation}
 \left\lfloor\frac{md}{h}\right\rfloor
 =\left\lfloor\frac{m(d+1)-1}{h}\right\rfloor.
 \label{act21:censo:prueba-estabilidad}
\end{equation}
When this holds, either side gives that value.
\end{lemma}
\begin{proof}
The minimum of \(mx\) is attained at the left endpoint. Its supremum
is \(m(d+1)/h\) and is not attained. The largest integer value that
\(\lfloor mx\rfloor\) can take is
\(\lceil m(d+1)/h\rceil-1\), which equals the second side because
the numerator and denominator are integers. Equality of the endpoints
of the range of values proves the assertion.
\end{proof}

The stability test holds in all three channels for
\(m=3^{30},Q,3^{1080}\). The last choice will be used solely
for a subsequent comparison, not to select the profile. The
first thirty recovered positions are precisely the five
six-trit stages of the regional transport:
\begin{center}\small
\begin{tabular}{@{}c l r@{}}\toprule
Channel&Thirty-trit prefix&\(P_\chi\)\\\midrule
\(\pi\)&\texttt{010211012222010211002111110221}&29152671743887\\
\(e\)&\texttt{201101121221102011012222102011}&147887858824447\\
\(\varphi\)&\texttt{121200112202121020010210010200}&127247717616687\\\bottomrule
\end{tabular}
\end{center}
The agreement preserves the prolongation already constructed; the decimal readout
is not used to replace its transport matrices retrospectively.

For a fixed channel, a tail of \(r\) trits is represented by
\(0\leq R<3^r\). Its cylinder and the cylinder of the first \(k\)
triads are
\begin{equation}
 I_{\chi,R}=\left[\frac{P_\chi3^r+R}{D},
                       \frac{P_\chi3^r+R+1}{D}\right),\qquad
 J_{\chi,T}=\left[\frac{T_\chi}{Q},\frac{T_\chi+1}{Q}\right).
 \label{act21:censo:dos-cilindros}
\end{equation}
The catalogue includes all tails for which these two cylinders
have a nonempty intersection.

\begin{proposition}[Catalogue by exact intersection]
\label{act21:censo:catalogos}
The compatible tails are the integers in the interval \([a_\chi,b_\chi]\),
where
\begin{align}
 a_\chi&=\max\left(0,
       \left\lfloor\frac{T_\chi D}{Q}\right\rfloor-P_\chi3^r\right),
       \label{act21:censo:extremo-a}\\
 b_\chi&=\min\left(3^r-1,
       \left\lfloor\frac{(T_\chi+1)D-1}{Q}\right\rfloor-P_\chi3^r\right).
       \label{act21:censo:extremo-b}
\end{align}
For the specified regional cylinders, no clipping applies and we obtain
\begin{equation}
\begin{array}{c|r|r|r}
 \chi&a_\chi\bmod729&b_\chi\bmod729&b_\chi-a_\chi+1\\\hline
 \pi&67&262&196\\
 e&584&51&197\\
 \varphi&117&312&196
\end{array}
\label{act21:censo:tabla-catalogos}
\end{equation}
\end{proposition}
\begin{proof}
Write \(N=P_\chi3^r+R\). The two half-open intervals in
\eqref{act21:censo:dos-cilindros} intersect exactly when
\[
 NQ<(T_\chi+1)D,\qquad (N+1)Q>T_\chi D.
\]
The second inequality is equivalent to
\(N\geq\lfloor T_\chi D/Q\rfloor\); the first, to
\(N\leq\lfloor((T_\chi+1)D-1)/Q\rfloor\). Subtracting
\(P_\chi3^r\) and retaining \(0\leq R<3^r\) proves the formulae.
Evaluation by integer division of the \(T_\chi\) extracted using
\eqref{act21:censo:prueba-estabilidad} gives the table. Since \(r\geq6\),
\(P_\chi3^r\equiv0\pmod{729}\), so the residue of the first
endpoint can also be obtained directly from
\(\lfloor T_\chi D/Q\rfloor\bmod729\).
\end{proof}

The last six trits of a tail form the word
\(\nu_6(R\bmod729)\), where \(\nu_6\) denotes the ternary representation
of length six, including its leading zeros. By the table, each
catalogue contains fewer than \(729\) consecutive elements, so
this reduction is injective on it. Thus, the three boundary catalogues
are completely described, without choosing the final triples:
\begin{equation}
 \begin{split}
 \mathcal B_\pi&=\{\nu_6(67+i):0\leq i<196\},\\
 \mathcal B_e&=\{\nu_6((584+j)\bmod729):0\leq j<197\},\\
 \mathcal B_\varphi&=\{\nu_6(117+\ell):0\leq\ell<196\}.
 \end{split}
 \label{act21:censo:fronteras-explicitas}
\end{equation}
Requiring the left endpoint of \(I_{\chi,R}\) to belong to
\(J_{\chi,T}\) would impose a different condition: it would lose the first cylinder
that intersects the boundary and would give \(195,196,195\). The census uses the
intersection of the complete cylinders, in accordance with its definition.

\subsubsection{Declared profile and compatibility indicators}

For \(u,v\in\mathbb F_3^6\), let
\(g(u,v)=\sum u_iv_i\in\mathbb F_3\),
\(n(u)=g(u,u)\), \(w(u)=\#\{i:u_i\ne0\}\), and
\(h(u,v)=\#\{i:u_i\ne v_i\}\). The last two functions take
integer values; the first two take values in \(\mathbb F_3\).
The reader's local profile, in the order \((\pi,e,\varphi)\), is
\begin{equation}
 \begin{split}
 (g_{\pi e},g_{\pi\varphi},g_{e\varphi})&=(2,2,0),\qquad
 (n_\pi,n_e,n_\varphi)=(0,0,0),\\
 (w_\pi,w_e,w_\varphi)&=(3,3,3),\qquad
 (h_{\pi e},h_{\pi\varphi},h_{e\varphi})=(2,5,3).
 \end{split}
 \label{act21:censo:perfil-local}
\end{equation}
The words are ordered as in \eqref{act21:censo:fronteras-explicitas}.
We use \([E]\in\{0,1\}\) for the indicator of a condition \(E\).
For each pair of channels \(\chi,\psi\), the rectangular matrix
\(A_{\chi\psi}\) has entry \([g(u,v)=g_{\chi\psi}]\).
The matrix \(A^{h}_{\chi\psi}\) multiplies that entry by
\([h(u,v)=h_{\chi\psi}]\). Let \(D_\chi^n\) be the diagonal matrix
of \([n(u)=0]\), and \(D_\chi^{nw}\) that of
\([n(u)=0][w(u)=3]\).

\Needspace{18\baselineskip}
\begin{samepage}
\begin{theorem}[Exhaustive enumeration by finite sums]
\label{act21:censo:teorema-conteo}
The product of the three catalogues contains \(7\,567\,952\) triples.
The successive conditions in \eqref{act21:censo:perfil-local} leave
\begin{equation}
 7\,567\,952\longrightarrow275\,794\longrightarrow6235
 \longrightarrow3127\longrightarrow331.
 \label{act21:censo:cadena-local}
\end{equation}
These four counts are, respectively,
\begin{align}
 N_g&=\operatorname{tr}(A_{\pi e}A_{e\varphi}A_{\varphi\pi}),
 \nonumber\\
 N_{gn}&=\operatorname{tr}(D_\pi^n A_{\pi e}
             D_e^n A_{e\varphi}D_\varphi^n A_{\varphi\pi}),
 \nonumber\\
 N_{gnw}&=\operatorname{tr}(D_\pi^{nw} A_{\pi e}
             D_e^{nw} A_{e\varphi}D_\varphi^{nw} A_{\varphi\pi}),
 \nonumber\\
 N_{gnwh}&=\operatorname{tr}(D_\pi^{nw} A^h_{\pi e}
             D_e^{nw} A^h_{e\varphi}D_\varphi^{nw} A^h_{\varphi\pi}).
 \label{act21:censo:trazas}
\end{align}
All products and traces in this formula are calculated over the integers,
not over the field with three elements.
\end{theorem}
\end{samepage}
\begin{proof}
The size of the product is \(196\cdot197\cdot196\). When a trace in
\eqref{act21:censo:trazas} is expanded, its indices run through exactly one
triple of words. The summand equals one if it satisfies all the represented
conditions and zero otherwise. Each triple is counted once.
The diagonal matrices add the individual conditions; the rectangular
factors add the pairwise conditions. This proves that the
traces are precisely the stated cardinalities.

To make their evaluation explicit, fix the indices \(i,j\) of the first
two channels and form the sets of indices \(\ell\) of the third
that satisfy each pairwise condition. Their indicators are intersected,
the relevant individual indicators are added, and the cardinality is summed.
The final sum runs over \(0\leq i<196\), \(0\leq j<197\).
Substitution of the words in \eqref{act21:censo:fronteras-explicitas} and
the profile \eqref{act21:censo:perfil-local} gives, in that order,
\(275794,6235,3127,331\). The supplement records the four contributions
of each of the \(196\) indices \(i\) and the complete list of \(331\) triples;
its program evaluates these sums with exact binary indicators and checks
each survivor by scalar evaluation. No sample or tolerance is involved.
\end{proof}

\subsubsection{Labelled incidence and orientation charge}

For a surviving triple, let \(B\in\mathbb F_3^{3\times6}\) be its
matrix of ordered rows. The second profile requires
\begin{equation}
 \operatorname{rank}_{\mathbb F_3}B=3,\qquad
 w_{\mathrm{col}}(B)=(1,1,2,2,3,0),\qquad
 q_{\partial}(B)=(1,2,2,1,2,0),
 \label{act21:censo:perfil-columnas}
\end{equation}
where \(w_{\mathrm{col}}\) counts the nonzero elements in each
column and \(q_{\partial}\) sums its entries modulo three. This is a profile
of labelled columns; permuting them without transporting the profile defines a different
reader. Denote by \(\mathcal F\) the set of \(331\) triples
from the preceding theorem.

\begin{proposition}[Reduction of the census to the terminal pair]
\label{act21:censo:dos-matrices}
The sums of indicators over \(\mathcal F\), successively imposing
rank, column weights, and column sums, are
\begin{equation}
 \sum_{B\in\mathcal F}[\operatorname{rank}B=3]=331,
 \qquad N_{\mathrm{rango,pesos}}=18,
 \qquad N_{\mathrm{rango,pesos,sumas}}=2.
 \label{act21:censo:conteo-columnas}
\end{equation}
With catalogue indices starting at one, the two final triples are
\((120,197,157)\) and \((129,197,148)\), and their matrices are
\[
 B_0=\begin{pmatrix}0&2&0&2&2&0\\0&0&1&2&2&0\\1&0&1&0&1&0\end{pmatrix},
 \qquad
 B_1=\begin{pmatrix}0&2&1&0&2&0\\0&0&1&2&2&0\\1&0&0&2&1&0\end{pmatrix}.
\]
\end{proposition}
\begin{proof}
For each triple in the already enumerated set, the rank is computed by elimination
over \(\mathbb F_3\), the nonzero elements of the six
columns are counted, and their entries are summed modulo three. The products of the indicators
in \eqref{act21:censo:perfil-columnas} give \eqref{act21:censo:conteo-columnas}.
Substituting the two final indices into
\eqref{act21:censo:fronteras-explicitas} gives the displayed matrices. The complete
list of survivors and the exact filters in the supplement allow
the elimination to be reproduced without presupposing these two matrices as the domain.
\end{proof}

The orientation \(\sigma=(1,1,-1)=(1,1,2)\in\mathbb F_3^3\)
selects the line
\(\langle\sigma\rangle=\{(0,0,0),(1,1,2),(2,2,1)\}\).
The row charges of \(B_0,B_1\) are \((0,2,0)\) and \((2,2,1)\);
only the latter belongs to this line. This recovers the last step
of the selector in Section~\ref{act21:censo:dos-matrices}, now preceded by
the construction and enumeration of its complete domain. The stable readout
of \(\lfloor3^{1080}x_\chi\rfloor\), performed after the filters,
gives the boundaries \texttt{021020}, \texttt{001220}, \texttt{100210},
with the same ranks \((129,197,148)\). This subsequent comparison does not define
the profile or enter the sums of indicators.

\subsubsection{Preserved information and scope of the result}

The procedure preserves the thirty-trit prefix, the compatible tail,
the regional order, the position of each trit, and the orientation. Reduction
modulo \(729\) preserves the identity of each candidate here because the
interval of tails has length less than \(729\); this injectivity
is not asserted for an arbitrary history. The census of \(331\) also proves
that the local profile alone does not select a single matrix: the
labelled marginals and the charge supply additional effective conditions.

Thus, a complete finite enumeration and its selector have been recovered
on output cylinders, with an explicit readout profile. Construction
of the signed register additionally uses its memory reader and its channels,
as specified below; the visible matrix is not identified with
that register by equality of the number of components. Nor does this count
prove the universality of the profile, generate the regional coordinates
anew, or decide a claim about alpha or about the complete continuum.
